\chapter{Methodology}
\label{chap:methodology}

This chapter documents the tools that produce every quantitative
result of the thesis.  \Cref{sec:bioverse} gives a high-level
overview of the \textsc{Bioverse} statistical framework
\parencite{BixelApai2021,Schlecker2024} that underpins the analysis
and highlights which parts this thesis inherits and which it extends.
The three subsequent sections walk through the pipeline in the same
order the data flows through it.  \Cref{sec:generator} details the
population \emph{generator}: the stellar sample, the assignment of
synthetic planets and orbits, the runaway greenhouse radius models,
the treatment of stellar ages, and the water-mass-fraction population
priors.  \Cref{sec:survey_simulator} covers the \emph{survey
simulator} that turns the underlying population into a mock observed
catalogue.  \Cref{sec:inference_methodology} finally documents the
two Bayesian tests used to quantify the \HZIED signal and the
multi-seed noise-realisation treatment that every sweep figure in
\Cref{chap:results} adopts.


% =======================================================================
\section{Bioverse Pipeline Overview}
\label{sec:bioverse}

The statistical framework of this thesis is built on \textsc{Bioverse},
the exoplanet-survey simulation pipeline of \textcite{BixelApai2021}
as extended for the \HZIED forecast by \textcite{Schlecker2024}.
Bioverse follows a strictly sequential three-stage architecture.  A
\emph{generator} draws a synthetic underlying population of stars
and planets from occurrence-rate priors and injects the physical
runaway-greenhouse radius inflation into the hot-side planets.  A
\emph{survey simulator} then projects that ``ground truth''
population into a mock observed catalogue by applying transit
selection and measurement noise.  Finally, a \emph{hypothesis-testing}
module computes the Bayesian evidence for competing demographic
models on the mock data via nested sampling.

The clean separation between the three stages has two operational
consequences that matter for the thesis.  First, the expensive step
is the generator: every call has to build a full stellar and
planetary catalogue and apply the physical radius-inflation
prescription.  The survey simulator and the inference module are
comparatively cheap.  Intermediate parquet-cached tables therefore
store the generated catalogue on disk between the generator and the
inference, so that repeated Bayesian runs at different injection
parameters or noise seeds do not incur the generator cost --- the
mechanism that makes the multi-seed sweeps of \Cref{chap:results}
computationally tractable.  Second, because the extensions this
thesis introduces all live within the generator step, the survey
simulator and the hypothesis-testing module can be inherited from
Schlecker+2024 essentially unchanged.

This thesis inherits the Bioverse framework as its baseline and
extends it in four directions, all within the generator:
\begin{enumerate}
    \item the generic Gaia sphere is replaced with the observed
        \PLATO P1--P5 target catalogues (\Cref{ssec:bioverse_stars};
        target selection is deferred to \Cref{chap:plato_mission});
    \item the uniform stellar-age prior is replaced with per-star
        ages derived from Gaia DR3 FLAME isochrones
        (\Cref{sec:age_modelling});
    \item four alternative water-mass-fraction population priors are
        introduced in place of the single fiducial injection value
        of Schlecker+2024 (\Cref{sec:composition_prior}); and
    \item the runaway-greenhouse radius model is upgraded with the
        \textcite{Aguichine2021} and Hycean extensions
        (\Cref{sec:runaway_models}).
\end{enumerate}
The survey simulator (\Cref{sec:survey_simulator}) and the
hypothesis-testing module (\Cref{sec:inference_methodology}) are
inherited from Schlecker+2024 unchanged.


% =======================================================================
\section{The Generator}
\label{sec:generator}

The generator produces an underlying stellar and planetary population
before any observational selection is applied.  It runs in five
substeps: build a stellar sample (\Cref{ssec:bioverse_stars}), draw
planets around each star from an occurrence-rate prior
(\Cref{ssec:bioverse_planets}), apply the runaway-greenhouse radius
inflation prescription (\Cref{sec:runaway_models}), attach stellar
ages and derive per-planet post-runaway ages
(\Cref{sec:age_modelling}), and draw per-planet water inventories
from a composition prior (\Cref{sec:composition_prior}).  Its output
is a table of ``ground truth'' planet properties on which the survey
simulator will later impose transit selection and radius noise.


\subsection{Stellar sample}
\label{ssec:bioverse_stars}

The baseline stellar sample of \textcite{Schlecker2024} is drawn from
the Gaia-based catalogue of \textcite{HardegreeUllman2023}, itself
built on the Gaia Catalogue of Nearby Stars
\parencite{Smart2021}, and covers a spherical volume of radius
$d_{\max}=120\,\mathrm{pc}$ around the Sun.  Effective temperatures,
luminosities, radii, and masses are derived from Gaia $G$,
$G_{\mathrm{BP}}$, $G_{\mathrm{RP}}$ photometry combined with 2MASS
$K_S$ magnitudes.  \textcite{HardegreeUllman2023} report that the
derived stellar parameters are consistent to within a few percent
with independently measured values for known exoplanet hosts, which
is what makes the catalogue suitable for a demographic study of
small-planet radii.

This thesis replaces the generic Gaia sphere with the observed
\PLATO Input Catalogue P1--P5 target lists documented in
\Cref{chap:plato_mission}.  The PLATO catalogues span a much wider
distance range than $120\,\mathrm{pc}$, and their target selection
is driven by asteroseismic characterisability rather than distance;
running Bioverse on them therefore requires disabling the $d_{\max}$
cut but leaves every other generator step untouched.


\subsection{Planetary population and orbits}
\label{ssec:bioverse_planets}

Each star is assigned a synthetic planetary system according to the
occurrence-rate priors of \textcite{Bergsten2022}, which parameterise
the per-star planet frequency as a function of planetary radius,
orbital period, and stellar spectral type from \Kepler statistics.
For every star Bioverse draws a per-system planet multiplicity, and
for every planet an orbital period and a baseline radius from this
parameterisation.  The semi-major axis follows from Kepler's third
law using the parent star's mass, and the instellation is computed
from the stellar luminosity via the inverse-square law.  Every
synthetic planet therefore inherits its full stellar-parameter
context, so that the host-dependent runaway threshold
$S_\mathrm{thresh}$ can be evaluated self-consistently in
\Cref{sec:runaway_models}.


\subsection{Runaway greenhouse radius models}
\label{sec:runaway_models}

The physical basis of the \HZIED signal is the Simpson--Nakajima
outgoing longwave radiation limit \parencite{Nakajima1992,Goldblatt2013}:
above a critical absorbed flux $S_\mathrm{thresh}$, no steady-state
radiative--convective equilibrium can sustain a saturated
water-vapour atmosphere in balance with the incoming stellar flux
while maintaining a liquid ocean at the surface
(\Cref{sec:habitable_zone}).  Planets crossing this threshold undergo
a runaway greenhouse and settle into a
steam-atmosphere-over-magma-ocean state whose transit radius is
inflated relative to the corresponding cold-side planet by tens of
percent.

The default radius-inflation prescription used throughout this
thesis follows \textcite{Turbet2020}, who computed a tabulated
mass--radius grid for water-rich rocky planets more irradiated than
the runaway greenhouse limit.  Each entry gives the inflated transit
radius as a function of planet mass and water mass fraction
$x_{\mathrm{H_2O}}$, spanning
$10^{-4}\leq x_{\mathrm{H_2O}}\leq 5\times 10^{-2}$.
\textcite{DornLichtenberg2021} added the constraint that a
substantial fraction of the outgassed water dissolves back into the
silicate melt during the magma-ocean phase and does not contribute
to the atmospheric column, reducing the effective radius inflation
by up to a factor of two at $x_{\mathrm{H_2O}}\sim 10^{-2}$.  The
combination --- Turbet+2020 with the Dorn--Lichtenberg
water-retention correction --- is what \textcite{Schlecker2024}
adopted, and it remains the fiducial baseline of this thesis.

Two extensions of the model space are also implemented and used in
the sensitivity analyses of \Cref{chap:results}.  For water-dominated
interior compositions ($x_{\mathrm{H_2O}}$ up to ${\sim}\,0.5$),
\textcite{Aguichine2021} extended the mass--radius grid to the
ocean-planet regime, providing an alternative prescription in which
the transit radius is set by a supercritical-water mantle rather
than by a steam atmosphere.  For hydrogen-dominated (``Hycean'')
atmospheres above a supercritical water interior,
\textcite{Pierrehumbert2022} and \textcite{Innes2023} show that the
runaway threshold shifts to a substantially different instellation
\openquestion{verify sign and value: Innes 2023 abstract states the
Hycean runaway threshold is \emph{reduced} (inner edge moves from
$\sim 1\,\mathrm{AU}$ to $1.6\,\mathrm{AU}$ for a G-star), whereas
the current text implies an increase to
$\sim 435\,\mathrm{W\,m^{-2}}$ --- verify against
\textcite{Pierrehumbert2022}}, and that the resulting radius
inflation profile is qualitatively different from the rocky case.
\Cref{chap:results} compares these three prescriptions systematically
on the same \PLATO forecast samples, and \Cref{ssec:hycean_mixed}
explores a mixed population in which the rocky and Hycean
prescriptions coexist.

The runaway-greenhouse fraction $f_\mathrm{rgh}$ enters the model at
this stage: for each hot-side planet, a Bernoulli draw with
probability $f_\mathrm{rgh}$ decides whether the tabulated inflated
radius is assigned in place of the baseline dry-rocky radius.
$f_\mathrm{rgh}$ is a lump parameter absorbing the population-level
uncertainty in whether a runaway planet retains its inflated steam
atmosphere on gigayear timescales; setting $f_\mathrm{rgh}=1$
represents the strong-signal limit in which every hot-side planet
retains its steam atmosphere, while smaller values dilute the
demographic step in proportion to the fraction of planets that have
lost their atmosphere on the survey timescale, allowing the \HZIED
test to remain meaningful even when atmospheric persistence is
uncertain.


% \subsection{Stellar ages and time-since-runaway} removed for condensed version
% \label{sec:age_modelling}  % kept as phantom for cross-references
\phantomsection\label{sec:age_modelling}

% \subsection{Water-mass-fraction population priors} removed for condensed version
% \label{sec:composition_prior}  % kept as phantom for cross-references
\phantomsection\label{sec:composition_prior}


% =======================================================================
\section{Survey Simulator}
\label{sec:survey_simulator}

The survey simulator converts the underlying population produced by
the generator into a mock observed catalogue --- the object on which
the hypothesis test will operate.  Three effects are applied in
sequence.

First, the geometric transit probability
$p_{\mathrm{tr}} \approx R_\star / a$ retains only those planets whose
orbital plane is close enough to the line of sight for a transit to
occur; for a habitable-zone orbit around a solar-analog host this
retains roughly $0.5\%$ of the underlying planets.  Second, an
SNR-limited detection cut removes planets whose transit signal,
integrated over the assumed mission duration, falls below the
detection threshold of the simulated survey; the noise floor is
parameterised as a per-star photometric standard deviation
consistent with the chosen mission architecture.  Third, a Gaussian
radius-measurement noise with per-planet standard deviation
$\sigma_R$ is added to the true radius; this is the dominant driver
of the individual-planet uncertainty budget for small, long-period
transiters and is what the hypothesis test in
\Cref{sec:inference_methodology} must marginalise over.

The three steps are chained in order and produce a table with the
same column schema as the generator output, but with fewer rows
(from the transit and SNR cuts) and a noise-broadened radius column.
That table is what the hypothesis-testing module takes as input.
The two selection effects intrinsic to the geometry --- the low
transit probability of long-period planets and the small transit
depth of Earth-sized bodies --- both work against the observation of
the \HZIED signature; the corresponding demographic biases are
discussed in \Cref{sec:transit_method} and quantified for the
\PLATO samples in \Cref{chap:results}.


% =======================================================================
\section{Hypothesis Testing}
\label{sec:inference_methodology}

Every number quoted in \Cref{chap:results} comes out of the same
top-level entry point, \texttt{run\_hzied\_analysis}, which supports
two hypothesis tests and a shared noise-realisation protocol.  The
subsection immediately below sets out why a Bayesian model
comparison is the natural framework for a demographic test of this
kind; the three subsections after it describe the tests themselves
and the multi-seed noise-realisation protocol used throughout the
results chapter.

% Subsections merged into one flat section for condensed version
% \subsection{Why Bayesian model comparison}
\phantomsection\label{ssec:why_bayesian}

\emph{Scaffold.}
Frequentist significance testing versus Bayesian model comparison for
population-level inferences.  The Bayesian evidence $Z$ as the
marginal likelihood of a model, and the log-evidence difference
between the step hypothesis $H_1$ and the flat null $H_0$,
%
\begin{equation}
    \ln\Delta Z = \ln \frac{\mathcal{Z}(H_1)}{\mathcal{Z}(H_0)},
    \label{eq:lnDZ}
\end{equation}
%
as the natural comparison statistic \parencite{Skilling2004}.  Nested
sampling via \texttt{dynesty} \parencite{Speagle2020} as the
computational engine, chosen for its robust treatment of degenerate
posteriors that plague the \HZIED parameter space at low
$x_{\mathrm{H_2O}}$.  Detection thresholds in terms of $\ln\Delta Z$:
moderate, strong, and decisive \parencite{Trotta2008}.
\textcite{Schlecker2024} established that a sample of
$N \gtrsim 100$ HZ planets with a few-percent radius precision and
$f_\mathrm{rgh} \gtrsim 0.1$ is sufficient for a decisive detection
($\ln\Delta Z \gtrsim 10$).  Why population-level tests are more
robust than individual-planet inferences for questions about
habitability: signal averaging over per-planet stochastic variation,
statistical power scaling with sample size, and the ability to place
informative upper limits in the case of a non-detection.


% \subsection{Full SMA-based test}
\phantomsection\label{ssec:sma_test}

The primary test in this thesis is the sliding-mean-average
construction of \textcite{Schlecker2024}.  Planets are sorted by
observed instellation $S_\mathrm{obs}$ and the observed radius
$R_\mathrm{obs}$ is smoothed with a rolling window of $w=25$
neighbours; \textsc{dynesty} \parencite{Speagle2020} then computes
the evidence ratio between the step model H\textsubscript{1} ---
with free parameters $S_\mathrm{thresh}$, $x_{\mathrm{H_2O}}$,
$f_\mathrm{rgh}$, and a baseline average radius $\bar R$ --- and the
flat null H\textsubscript{0} which shares only $\bar R$.  All priors
are uniform except for $x_{\mathrm{H_2O}}$, which is log-uniform
over the AGNI tabulated range, and every run uses
$n_\mathrm{live}=100$.

The rolling smoother makes the test powerful because it uses the
full shape of the radius profile on both sides of the threshold.
Its weakness is that adjacent window positions overlap by
construction, so the Gaussian likelihood assumes independent data
that isn't actually independent; the resulting $\ln\Delta Z$
therefore overstates the evidence relative to a strict likelihood,
and needs the sanity check described next.


% \subsection{Simple two-mean test at fixed $S_\mathrm{thresh}$}
\phantomsection\label{ssec:hypothesis_test_methodology}

Sometimes we want a Bayes factor that treats $S_\mathrm{thresh}$ as
known rather than fitting for it --- for a rough diagnostic, or to
test the SMA result against the correlation-free lower bound.
Setting \texttt{hypothesis\_mode='means'} in
\texttt{run\_hzied\_analysis} switches from the full SMA inference
to a two-aggregate Gaussian mean-comparison.

The sample is split at the injected $S_\mathrm{thresh}$ into a cold
subset and a hot subset.  Each side contributes exactly one aggregate
data point --- its mean observed radius and the standard error of
that mean.  The comparison then pits H\textsubscript{1} with two free
means ($\mu_\mathrm{cold}$, $\mu_\mathrm{hot}$) against
H\textsubscript{0} with a single free mean $\mu$.  Priors on each free
mean are uniform on $[0.1,\,15]\,\Rearth$, and the Bayes factor comes
either from a dynesty run on the two data points or from the
closed-form Gaussian expression appropriate when the priors are
broad.  The two agree to within sampler noise, and the closed-form
version is what \Cref{chap:kepler_tess} uses on the real Kepler/TESS
sample because it does not require the pipeline.

By construction the two-mean test throws away the profile shape and
keeps only two aggregates, so its Bayes factor is a lower bound on
what the SMA test can extract.  It also acts as a check on the SMA
number: an SMA $\ln\Delta Z$ that looks decisive but corresponds to
an under-powered two-mean test is a sign that the rolling smoother
is inflating the evidence rather than picking up a real step.


% \subsection{Recovering $\hat{S}_\mathrm{thresh}$ from the posterior}
\phantomsection\label{ssec:sthresh_recovery}

The SMA chains carry $S_\mathrm{thresh}$ as their first parameter,
and every sweep figure in \Cref{chap:results} with a right-hand
panel simply reports the posterior median of that column per noise
realisation.  Comparing that recovered $\hat{S}_\mathrm{thresh}$
against the injected value gives a second diagnostic that is largely
orthogonal to $\ln\Delta Z$: if the evidence is decisive and the
threshold recovers near the injected value, the detection is
trustworthy; if the evidence is decisive but $\hat{S}_\mathrm{thresh}$
drifts wildly, the single-step model is compensating for structure
it cannot represent.  \Cref{ssec:hycean_mixed} shows a concrete
example in the high-$f_\mathrm{hycean}$ regime, where $\ln\Delta Z$
climbs past $200$ but the recovered threshold overshoots the
injected $280\,\mathrm{W\,m^{-2}}$ by hundreds of units.


% \subsection{Multi-seed noise-realisation treatment}
\phantomsection\label{ssec:multi_seed}

One call to the pipeline uses one measurement-noise draw and one
sampler seed, and the resulting $\ln\Delta Z$ is more variable than
one might guess.  On the wrr sweep of \Cref{ssec:wrr_sweep}, for
instance, changing the seed at fixed injection shifts the reported
Bayes factor by roughly the same amount as changing the injected
water content by half an order of magnitude.  Reading any
single-seed number as a definitive value would therefore be
misleading.

Every sweep figure in \Cref{chap:results} runs ten independent noise
realisations per grid point.  The loop picks a per-grid-point seed
base $s_i = s_0 + \Delta\,i$ (with $\Delta=200$, generous enough
that no two grid points can collide in the inference cache even
when the planet count happens to be constant across the sweep ---
as on P4 throughout the wrr sweep), then runs \texttt{n\_trials=10},
which internally uses seeds $s_i + 13\,k$ for $k = 0,\ldots,9$.  The
resulting ten $\ln\Delta Z$ and ten $\hat{S}_\mathrm{thresh}$ values
are collapsed to a median and a $16$--$84\,\%$ percentile band per
grid point, and those are what the sweep figures plot.

Every Bayes factor quoted for a specific point in this thesis is
therefore a median across ten seeds, and the accompanying
seed-to-seed range should be read as the effective uncertainty at
$n_\mathrm{live}=100$.


% =======================================================================
\section{Simulation Configuration}
\label{sec:simulation_config}

\emph{Scaffold.}
Fiducial choices adopted throughout this work: \PLATO P4 as the
headline sample, P5 as the secondary sample;
\texttt{agni\_updated} radius inflation with the
\textcite{DornLichtenberg2021} water-incorporation correction;
strict FLAME age assignment.  Per-sample radius precision drawn from
the \emph{\PLATO Definition Study Report} \parencite{ESA2017PLATO}.
Observable cuts on planet mass, transit depth, and instellation.
Complete tabulation of the fiducial parameter set.
